\section*{Chapitre \ref{ch:g10:electron-shells} --- Les couches électroniques et le tableau périodique}

\begin{solution}{exo:g10:electron-shells:1}
C~: $1\mathrm{s}^2\,2\mathrm{s}^2\,2\mathrm{p}^2$. N~:
$1\mathrm{s}^2\,2\mathrm{s}^2\,2\mathrm{p}^3$. Mg~:
$1\mathrm{s}^2\,2\mathrm{s}^2\,2\mathrm{p}^6\,3\mathrm{s}^2$.
\end{solution}

\begin{solution}{exo:g10:electron-shells:2}
Carbone 4, azote 5, magnésium 2.
\end{solution}

\begin{solution}{exo:g10:electron-shells:3}
\omterm{def:g10:electron-shells:configuration}{Couche} externe $n = 3$~: \omterm{def:g9:periodic-table-first-look:period}{période} 3. $2 + 3 = 5$ \omterm{def:g10:electron-shells:valence}{électrons de valence}~:
\omterm{def:g9:periodic-table-first-look:period}{groupe} 15 (le phosphore).
\end{solution}

\begin{solution}{exo:g10:electron-shells:4}
Les \omterm{def:g10:electron-shells:family}{métaux alcalins} (lithium, sodium, potassium)~; les \omterm{def:g10:electron-shells:family}{halogènes}
(fluor, chlore, brome, iode)~; les \omterm{def:g10:electron-shells:family}{gaz nobles} (hélium, néon, argon).
\end{solution}

\begin{solution}{exo:g10:electron-shells:5}
Les \omterm{def:g7:atoms-and-molecules:atom}{atomes} tendent à gagner ou à perdre des \omterm{def:g9:inside-the-atom:nucleus}{électrons} de façon à avoir
huit \omterm{def:g9:inside-the-atom:nucleus}{électrons} sur leur \omterm{def:g10:electron-shells:configuration}{couche} externe, comme le \omterm{def:g10:electron-shells:family}{gaz noble} le plus
proche.
\end{solution}

\begin{solution}{exo:g10:electron-shells:6}
K, $2.8.8.1$, en perd un~: \ce{K+},
$1\mathrm{s}^2\,2\mathrm{s}^2\,2\mathrm{p}^6\,3\mathrm{s}^2\,3\mathrm{p}^6$
(argon). F, $2.7$, en gagne un~: \ce{F-},
$1\mathrm{s}^2\,2\mathrm{s}^2\,2\mathrm{p}^6$ (néon).
\end{solution}

\begin{solution}{exo:g10:electron-shells:7}
$1\mathrm{s}^2\,2\mathrm{s}^2\,2\mathrm{p}^6\,3\mathrm{s}^2$~: $Z = 12$,
le magnésium.
\end{solution}

\begin{solution}{exo:g10:electron-shells:8}
\ce{Mg^{2+}} et \ce{Cl-}~: \ce{MgCl2}. \ce{Al^{3+}} et \ce{O^{2-}}~:
\ce{Al2O3}.
\end{solution}

\begin{solution}{exo:g10:electron-shells:9}
Le soufre (6 \omterm{def:g10:electron-shells:valence}{électrons de valence}, $2.8.6$), dans le même groupe, le 16.
\end{solution}

\begin{solution}{exo:g10:electron-shells:10}
Sa \omterm{def:g10:electron-shells:configuration}{couche} externe est déjà complète (huit \omterm{def:g9:inside-the-atom:nucleus}{électrons})~: il n'a rien à
gagner à perdre ou à gagner des \omterm{def:g9:inside-the-atom:nucleus}{électrons}.
\end{solution}

\begin{solution}{exo:g10:electron-shells:11}
L'argon a 18 \omterm{def:g9:inside-the-atom:nucleus}{électrons}~; \ce{X^{2+}} en a perdu 2, donc X en a 20~: c'est
le calcium.
\end{solution}

\begin{solution}{exo:g10:electron-shells:12}
L'hélium, $1\mathrm{s}^2$, a une \omterm{def:g10:electron-shells:configuration}{couche} externe complète (la \omterm{def:g10:electron-shells:configuration}{couche} 1
ne contient que deux \omterm{def:g9:inside-the-atom:nucleus}{électrons}), comme les autres \omterm{def:g10:electron-shells:family}{gaz nobles}~; il ne se
comporte pas comme les métaux du \omterm{def:g9:periodic-table-first-look:period}{groupe} 2.
\end{solution}

\begin{solution}{exo:g10:electron-shells:13}
Son \omterm{def:g10:electron-shells:valence}{électron de valence} est dans la \omterm{def:g10:electron-shells:configuration}{couche} 4, plus loin du \omterm{def:g9:inside-the-atom:nucleus}{noyau} que
celui du sodium, dans la \omterm{def:g10:electron-shells:configuration}{couche} 3~: moins retenu, il est perdu plus
facilement.
\end{solution}

\begin{solution}{exo:g10:electron-shells:14}
Les quatre ont 18 \omterm{def:g9:inside-the-atom:nucleus}{électrons} et la configuration de l'argon~:
\[ 1\mathrm{s}^2\,2\mathrm{s}^2\,2\mathrm{p}^6\,3\mathrm{s}^2\,3\mathrm{p}^6. \]
\end{solution}

\begin{solution}{exo:g10:electron-shells:15}
Il gagne 3 \omterm{def:g9:inside-the-atom:nucleus}{électrons} pour atteindre un octet~: il a 5 \omterm{def:g10:electron-shells:valence}{électrons de
valence}, \omterm{def:g9:periodic-table-first-look:period}{groupe} 15. Dans la \omterm{def:g9:periodic-table-first-look:period}{période} 3, c'est le phosphore~:
\[ 1\mathrm{s}^2\,2\mathrm{s}^2\,2\mathrm{p}^6\,3\mathrm{s}^2\,3\mathrm{p}^3. \]
\end{solution}

\begin{solution}{pb:g10:electron-shells:1}
\textbf{1.} Aluminium~: 13 \omterm{def:g9:inside-the-atom:proton}{protons}, 13 \omterm{def:g9:inside-the-atom:nucleus}{électrons}. Oxygène~: 8 \omterm{def:g9:inside-the-atom:proton}{protons}, 8
\omterm{def:g9:inside-the-atom:nucleus}{électrons}.

\textbf{2.} Al~: $1\mathrm{s}^2\,2\mathrm{s}^2\,2\mathrm{p}^6\,3\mathrm{s}^2\,3\mathrm{p}^1$.
O~: $1\mathrm{s}^2\,2\mathrm{s}^2\,2\mathrm{p}^4$.

\textbf{3.} Al~: 3 \omterm{def:g10:electron-shells:valence}{électrons de valence}, \omterm{def:g9:periodic-table-first-look:period}{période} 3, \omterm{def:g9:periodic-table-first-look:period}{groupe} 13. O~: 6,
\omterm{def:g9:periodic-table-first-look:period}{période} 2, \omterm{def:g9:periodic-table-first-look:period}{groupe} 16.

\textbf{4.} L'aluminium est un \omterm{def:g9:periodic-table-first-look:metal}{métal}, l'oxygène un \omterm{def:g9:periodic-table-first-look:metal}{non-métal}.

\textbf{5.} \ce{Al^{3+}}~: $1\mathrm{s}^2\,2\mathrm{s}^2\,2\mathrm{p}^6$.

\textbf{6.} \ce{O^{2-}}~: $1\mathrm{s}^2\,2\mathrm{s}^2\,2\mathrm{p}^6$.

\textbf{7.} Le néon.

\textbf{8.} L'aluminium en perd 3, l'oxygène en gagne 2.

\textbf{9.} $2 \times (+3) + 3 \times (-2) = 0$~: \ce{Al2O3}.

\textbf{10.} Perdus~: $2 \times 3 = 6$. Gagnés~: $3 \times 2 = 6$.

\textbf{11.} Les \omterm{def:g9:inside-the-atom:nucleus}{électrons} ne sont ni créés ni détruits~: ceux que perd
l'aluminium sont exactement ceux que gagne l'oxygène, et le composé est
neutre.

\textbf{12.} Il porte la même charge, $+3$~: les charges du cristal
restent équilibrées.

\textbf{13.} $2 \times 27 + 3 \times 16 = 102$ \omterm{def:g9:inside-the-atom:proton}{nucléons}.

\textbf{14.} $102 \times 1.67 \times 10^{-27} = \qty{1.70e-25}{kg}$.

\textbf{15.} $10^{-3} / 1.70 \times 10^{-25} \approx 5.87 \times 10^{21}$
unités de formule.

\textbf{16.} $2 \times 5.87 \times 10^{21} = 1.17 \times 10^{22}$ \omterm{def:g9:ions:ion}{ions}
\ce{Al^{3+}}~; $3 \times 5.87 \times 10^{21} = 1.76 \times 10^{22}$ \omterm{def:g9:ions:ion}{ions}
\ce{O^{2-}}.

\textbf{17.} Un \omterm{def:g9:inside-the-atom:nucleus}{électron} est environ 1836 fois plus léger qu'un \omterm{def:g9:inside-the-atom:proton}{nucléon}~;
les 50 \omterm{def:g9:inside-the-atom:nucleus}{électrons} d'une unité de formule pèsent moins de trois centièmes
d'un \omterm{def:g9:inside-the-atom:proton}{nucléon}.

\textbf{18.} $6 \times 5.87 \times 10^{21} \approx 3.5 \times 10^{22}$
\omterm{def:g9:inside-the-atom:nucleus}{électrons} transférés pour \qty{1}{g} de corindon.
\end{solution}
